%==========================================================
\section{Uso de IA generativa}\label{anx:ia}

O enunciado pede que, se a equipa usou um LLM para ajudar a agrupar notas de campo ou a redigir o guião do percurso cognitivo, mostre \textbf{um exemplo do que a IA errou}. Este anexo regista onde a IA foi usada, o que foi verificado à mão e o que teria ficado errado sem essa verificação.

\subsection{Onde foi usada}

\begin{table}[H]
\centering\small
\begin{tabularx}{\textwidth}{@{}p{3.4cm}LL@{}}
\toprule
\textbf{Tarefa} & \textbf{O que a IA fez} & \textbf{O que a equipa verificou ou mudou} \\
\midrule
Revisão dos registos de entrevista & Detetou incoerências entre registos, notas de interpretação e siglas (referências a registos inexistentes, siglas sem definição). & Cada correção foi aceite ou rejeitada pela equipa; a menção ao ChatGPT foi mantida contra a sugestão da IA. \\
Modelos culturais & Redesenhou os modelos com a notação do livro (direção das setas, voz do influenciador, \textit{pushback}, sobreposição). & Cada influência foi conferida com o registo que cita. \\
Consolidação do fluxo e do cultural & Propôs papéis, agrupamentos de influenciadores e nomes abstratos para comunicações e artefactos. & Papéis e agrupamentos revistos contra os modelos individuais. \\
Modelo do relatório & Gerou a estrutura em \LaTeX{} e o texto das secções já com dados. & \textit{a preencher} \\
Diagrama de afinidade & \textit{a preencher, se usada} & \\
Avaliação heurística e percurso cognitivo & \textit{a preencher, se usada} & \\
\bottomrule
\end{tabularx}
\caption{Uso de IA generativa nesta entrega.}
\label{tab:uso-ia}
\end{table}

\subsection{Exemplo do que a IA errou}

\pendente{Equipa}{%
Escolher o exemplo a apresentar e completá-lo (o que a IA produziu, porque estava errado, como foi detetado e como foi corrigido). Casos registados durante esta entrega:
\begin{itemize}
  \item \textbf{Sigla inventada.} O documento de A0 usava «F/OP» sem definir o «F». A IA propôs «fala geral» e escreveu essa definição no glossário, sem base nos dados. A definição tem de ser confirmada ou corrigida pela equipa.
  \item \textbf{Classificação de registos} (só serve se essa versão tiver sido gerada com IA; confirmar). Na versão inicial de U0, acontecimentos passados noutros locais (por exemplo, «no corredor, entre aulas») estavam marcados como O (observado), embora o entrevistador não pudesse tê-los visto. Passaram a R (relato).
  \item \textbf{Contagem errada.} Ao resumir as sessões, a IA indicou 7 registos observados e 12 relatos em U0; a contagem correta é 8 e 11.
  \item \textbf{Pressuposto sobre a equipa.} Ao distribuir os papéis das sessões de interpretação, a IA assumiu quem tinha entrevistado A0 e U0 sem que isso estivesse nos dados.
\end{itemize}
Se a IA for usada no diagrama de afinidade, a comparação entre os agrupamentos da IA e os da equipa (padrões inventados, perda da origem das notas) é o exemplo mais relevante para o enunciado.}
